%DS0



\section{ĐỀ SỐ \thesection }
\setcounter{bt}{0}
%%Câu 1.
\begin{bt}  Thực hiện phép tính sau:
	
	\begin{listEX}[3]
		\item $-256 + 32 - 744$;
		\item $24\cdot15 + 24\cdot84 + 24$;
		\item $4\cdot5^2 - 3\cdot8 + 16:4 + 2017^0$.	
	\end{listEX}
	
	\loigiai{
		\begin{listEX}[3]
			\item $-256 + 32 - 744 \\= (-256 -744) + 32 \\= -1000 + 32  = -968$;
			\item $24\cdot15 + 24\cdot84 + 24\\ = 24 \cdot(15 + 84 + 1) \\= 2400$;
			\item  $4\cdot5^2 - 3\cdot8 + 16:4 + 2017^0 \\= 4\cdot25 - 24 + 4 + 1 \\= 81$.
		\end{listEX}
	}
\end{bt}


%%Câu 2.
\begin{bt} Tìm số nguyên $x$, biết
	\begin{listEX}[3]
		\item $4\cdot(x + 2) = 24$;
		\item $-x + 20 = 8$;
		\item $x \, \vdots \,  5$ và $9 < x < 30$.
	\end{listEX}
	
	\loigiai{
	\begin{listEX}[3]
			\item $4\cdot(x + 2)  = 24$\\
				\hspace*{0.85 cm}$\begin{aligned} 
				x + 2 &= 6\\
				x &= 4
				\end{aligned}$ 
			\item $-x + 20  = 8$\\
				\hspace*{1.25 cm}$\begin{aligned}
				x &= 20-8\\
				x &= 12
				\end{aligned}$
			\item 
				$x\,\vdots\,5 \Rightarrow x \in B(5)$\\
				$x \in \{10;15;20;25\}$
			\end{listEX}
	}
\end{bt}



%%Câu 3.
\begin{bt} Mai đi nhà sách mua dụng cụ học tập. Mai mua $5$ cuốn tập, $8$ cây viết bi, $2$ hộp bút màu. Biết giá một cuốn tập là $14\,000$ đồng, giá một hộp màu là $35\,000$ đồng, giá một cây viết bi là $7\,000$ đồng. Mai đưa cho quầy thu ngân $200\,000$  đồng. Hỏi Mai được trả lại bao nhiêu tiền?
	
	\loigiai{
		Tổng số tiền Mai đã mua hết là: $5\cdot14\,000 + 8\cdot7\,000 + 2\cdot 35\,000 = 196\,000$ (đồng) 
		\\Số tiền Mai được trả lại là: $200\,000 - 196\,000 = 4000$ (đồng)
	}
\end{bt}


%%Câu 4.
\begin{bt}
	Một trường tổ chức cho học sinh đi tham quan bằng ô tô, biết rằng nếu xếp $40$ người hay $45$ người vào một xe thì đều không dư. Tính số học sinh đi tham quan, biết số học sinh khoảng từ $600$ đến $800$.
	\loigiai{
		Gọi $x$ $(x\in  \mathbb{N})$ là số học sinh đi tham quan. Ta có $x \in BC(40;45)$ và $600 \leqslant x \leqslant 800)$\\
	Phân tích các số ra thừa số nguyên tố  $40=2^3\cdot5 \hspace{1cm}45=3^3\cdot5$
		\\ BCNN$(40;45) =2^3\cdot3^3\cdot5= 360$
		\\BC$(40;45) =$ B$(360)=\{0; 360;720;1080;... \}$
		\\Mà $x \in N, 600 \leqslant x < 800$
		\\Nên $x = 720$\\
		Vậy có 720 học sinh đi tham quan. 
	}
\end{bt}


%%Câu 5.
\begin{bt}Nhà bạn Hà có một mảnh vườn hình chữ nhật với chiều dài là $25$ m và chiều rộng là $18$ m. Bố bạn Hà muốn đào một cái ao hình vuông với cạnh là $6$ m ở giữa mảnh vườn để nuôi cá, phần còn lại để mẹ bạn trồng rau.
	\begin{enumerate}
		\item  Tính diện tích ao nuôi cá nhà Hà?
		\item Mẹ bạn Hà đi mua hạt giống để trồng rau, biết rằng cứ $1$ mét vuông đất gieo được $3$ gam hạt giống. Hỏi cần ít nhất bao nhiêu gam hạt giống để gieo hết phần đất trồng rau? 
	\end{enumerate}
	
	\loigiai{
		\begin{enumerate}
			\item Diện tích ao nuôi cá nhà bạn Hà: $6\cdot6 = 36$ (m$^2$)
			\item Diện tích mảnh vườn nhà bạn Hà là: $25\cdot18 = 450$ (m$^2$)
			\\ Diện tích phần đất trồng rau là: $450 - 36 = 414$ (m$^2$)
			\\ Lượng hạt giống cần để gieo hết phần đất trồng rau là: $414\cdot3 = 1242$ (gam)
		\end{enumerate}
	}
\end{bt}


%%Câu 6.
\begin{bt}Quan sát biểu đồ tranh dưới đây và trả lời câu hỏi: 
	\begin{center}
		\def\icon{\faApple* \ \ }
		\def\hicon{\tikz{\clip(-0.18,-0.19)rectangle(0.02,0.45);\path(0,0)node{\faApple*}}}
		\def\vetranh#1{\foreach \dem in{1,...,#1}{\icon \ }}{
			\begin{tabular}{|c|p{10cm}<{}|}\hline 
				\multicolumn{2}{|c|}{\bf Khối lượng táo một cửa hàng bán được của các ngày trong tuần}\\ \hline 
				\bf Ngày &\multicolumn{1}{c|}{\bf Số ki-lô-gam táo bán được}\\ \hline 
				Thứ 2 &\vetranh{6}\\ \hline 
				Thứ 3 &\vetranh{3} \hicon \\ \hline 
				Thứ 4 &\vetranh{5} \\ \hline 
				Thứ 5 &\vetranh{4} \hicon \\ \hline 
				Thứ 6 &\vetranh{4} \\ \hline 
				Thứ 7 &\vetranh{8}\\ \hline 
			\end{tabular}\\
			(Cho biết: \icon = 20 kg ;    \quad    \hicon \ \ = 10 kg)}
	\end{center}
	\begin{enumerate}
		\item Ngày nào cửa hàng bán được khối lượng táo nhiều nhất?
		\item Ngày nào cửa hàng bán được khối lượng táo ít nhất?
		\item Tính tổng khối lượng táo cửa hàng bán được trong tuần đó?
	\end{enumerate}
	
	\loigiai{
		\begin{enumerate}
			\item Ngày thứ 7 cửa hàng bán được khối lượng táo nhiều nhất.
			\item Ngày thứ 3 cửa hàng bán được khối lượng táo ít nhất.
			\item Khối lượng táo cửa hàng bán được trong tuần đó là: $30\cdot20 + 2\cdot10 = 620$ (kg)
		\end{enumerate}
	}
\end{bt}


%%Câu 7.
\begin{bt} Cho $A = 1 + 2 + 2^2 + 2^3 + 2^4 +...+ 2^{99} $. Chứng minh rằng $A$ không chia hết cho $7$. 
	\loigiai{
		$A = 1+2+2^2+2^3+2^4+...+2^{99}$
		\\$A = (1+2+2^2)+2^3(1+2+2^2)+...+2^{96}(1+2+2^2)+2^{97}$
		\\ $A = 7 + 2^3\cdot 7+...+2^{96}\cdot7+2^{97}$
		\\ $A = 7(1+2^3+...+2^{96})+2^{97}$ 
		\\Mà $7(1+2^3+...+2^{96})\vdots7$ và $2^{97}\not\vdots7$ 
	nên $A\, \not\vdots\, 7$
	}
\end{bt}




